\Needspace{10\baselineskip}
\centerline{\textsc{Ceres and Pallas.}}

\medskip
Monatliche Correspondenz. Vol. IX, pp.\ 246—250. March 1804.

\medskip
On p.\ 370 of last October’s issue, we already published an ephemeris of the geocentric path of \emph{Ceres} for 1804, calculated by Dr.\ \textsc{Gauss} from his elements for its orbit, refined for the ninth time. On p.\ 535 of the December issue, we promised a chart of this course by Inspector \textsc{Harding} in Lilienthal. The chart now accompanies this issue.

{\rightskip=0pt plus 1em
Dr.\ \textsc{Gauss} has since corrected the elements of the orbit of \emph{Pallas} once more using the latest observations. From these orbital elements, now refined for the seventh time, he calculated the ephemeris below. We will likewise follow it with a chart of the planet’s geocentric path drawn up by Inspector \textsc{Harding}. [The ephemeris runs from 30 April 1804 to 19 January 1805 and gives the right ascension, declination, distance from Earth, and brightness in the form used up to now.]\par}


At some points the difference from the ephemeris given by Prof.\ \textsc{Bode} on p.\ 91 of the Berlin Astronomical Yearbook for 1806 is considerable. Observing \emph{Pallas} will be somewhat easier again this year than it was last year. On 23 July this year it reaches the peak brightness of 20 June last year, and on 7 September it becomes more than $\frac{1}{4}$ brighter still. In addition, the nights are completely dark, and \emph{Pegasus} and \emph{Aquarius} are less crowded with stars than the region the planet occupied last year. This time it reached as early as 5 June the brightness at which it was first found last year, and on 18 May it reached the brightness of 10 October 1803, when Dr.\ \textsc{Olbers} last observed it.

The latest elements (VII), from which Dr.\ \textsc{Gauss} calculated the above ephemeris for the orbit of \emph{Pallas}, are as follows:

\begin{center}
\begingroup\fontsize{9}{11}\selectfont\setlength{\tabcolsep}{3pt}
\begin{tabular}{lr}
Epoch (Seeberg) 1803 . . . . . . & $221^\circ\ 29'\ 32''0$ \\
Daily tropical motion . . . & $770''0446$ \\
Annual tropical motion . . . & $78^\circ\ 4'\ 26''3$ \\
Aphelion 1803 . . . . . . . & $301\ 17\ 34.4$ \\
Eccentricity . . . . . . . . . & $0.2457396$ \\
Logarithm of the semimajor axis  & $0.4423790$ \\
Ascending node 1803 . . . . . . & $172^\circ\ 28'\ 13''7$ \\
Inclination of the orbit . . . . . . . & $34\ 38\ 1.1$ \\
\end{tabular}\endgroup
\end{center}
\SourceLeaf{252}%

In last December’s issue, on p.\ 536, we published a short list, drawn from \textsc{Piazzi}’s large star catalogue, of stars that will lie on or near the path of \emph{Ceres}. Here we likewise give a short selection from that catalogue of stars that will come very close to \emph{Pallas} and can be compared with it.

\begin{flushright}
[\textsc{von Zach}.]
\end{flushright}

\bigskip
